%==============================================================================
% 九、模型评价与推广
%==============================================================================
\section{模型评价与推广}

\subsection{模型优点}

第一，数据核验做得细。本文把题面与实测的六处差异，包括样本数、sheet 数、相关系数类型、交叉矩阵规模、五个理化列的量纲与标签尺度逐条核对并列表，交叉校验达到 $10^{-17}$ 量级，因此后续结论建立在可复核的数据之上，而不是建立在题面文字的转述之上。

第二，三问共用一条证据链。问题一的预测器不只用于报精度，它的失效诊断，即地板效应占比与 $R^2$ 的单调关系解释了问题二为什么有一半记录做不了化学判断；问题二的 Top-3 并集直接构成问题三的 75 维决策空间；问题三的代理又沿用同一套感官定义。三问的结论因此可以互相对照，不会出现各训一个模型、彼此矛盾的情况。

第三，结论都标了可信范围。问题一报区间而非点估计，问题二把一半的 Top-3 记录标为无法做化学判断，问题三以“代理误差大于前沿点间距”为由拒绝给出单点最优。这些否定性结论不是退让，而是本题样本量下的正确表述。

第四，精确解优先。问题三用混合整数线性规划求得的前沿支配了 92.17\% 的 NSGA-II 点、反向支配为 0，且精确解的非支配性逐点验证通过，说明在可以线性化的结构上不应退而用启发式算法。但超体积的对比要谨慎：精确侧只有 7 个点、启发式侧 115 个点，而超体积随点数增加而增大，因此 757.78 对 299.12 不能读成"优 2.5 倍"；真正可比的是支配率与逐点水平。

\subsection{模型不足}

第一，标签不能精确还原。附件4 的连续强度与 0--6 整数分之间不存在可复现的对应关系，本文只能给出近似的分位映射。实测每列最高频取值仅占 1.8\%、等于列最小值仅占 0.88\%，说明标签已被中心化并加过扰动，连“大量样品取最低档”这一结构都无法从数据直接验证——等级比例只能取自感官评价文献而非拟合。即使如此，精确同档率中位也只有 0.376，且附件4 没有整数真值，任何 kappa 都只是映射内的往返一致性，不能称为对面板真值的精度。

第二，问题一的精度与文献持平，但分布高度右偏。均值 0.2223 与文献\cite{schreurs2024}报告的 0.21（梯度提升）、0.22（最优树模型）相当，而本文用的是嵌套交叉验证、选型不进评估折，同等数值下口径更保守。但中位只有 0.1469，比均值低 0.0754：50 项中仅 6 项 $R^2>0.5$（酸度 0.7149、酒精 0.7108、乳酸 0.6768、乙酸 0.6326、苦味 0.5862、酒体 0.5328），有 19 项低于 0.10。按类别拆分，15 个嗅觉项均值 0.1490、中位 0.1194，其余 35 个味觉与口感项均值 0.2536。文献只报一个整体均值，看不出这一分裂；本文并列中位与均值，正是为了把“近半数感官项无法由化学谱解释”显式化。整体喜好项的 $R^2$ 仅 0.1098，仍属最难预测的一类。此外 50 项中只有 7 项的残差通过正态性检验，多数维度的误差结构并非同方差，限制了区间估计的适用面。

第三，问题二的归因不稳定。跨模型 Jaccard 中位仅 0.0974、最小 0.0175、最大也只有 0.2533，说明 Top-3 名单高度依赖模型族；150 条记录中 83 条因阈值表缺项无法做 OAV 检验。因此本文只能报告 19 条双条件交集，这不足以支撑一份完整的“关键化合物清单”。

第四，问题三未经实验验证。代理的逐项 $R^2$ 中位 0.1386、最大误差 3.1525，折合标准差 2.64 倍，且优化中不含化合物间的协同与拮抗，生成的化学谱只是配方调整的方向性候选，不能替代中试与感官评价。另外，五类缺陷物的约束在本文数据下全部退化为 $x_j\leq x^0_j$，敏感性分析显示去掉它们后没有一类会恶化、$f_1$ 代价为 0，即这组约束在本例中不活跃，只能防止更大改动幅度下的触碰；而 trans-2-壬烯醛的阈值只能引用文献，其单位按附件8 标为 mg/L，与样本取值范围不匹配，使该缺陷物只能保证“不恶化”而非“不超标”。

\subsection{模型推广}

本文的做法对同类“高维成分数据 + 主观感官标签 + 小样本”的问题有可迁移性。可迁移的是四个环节：一是用官方对照表核对对数变换与缺失机制，避免自行推断；二是用嵌套交叉验证把模型选择与评估分开，这是小样本下报告 $R^2$ 的诚实做法；三是在强共线下用重采样的入选频率稳定化逐列置换，阈值由稳定选择的界反解；四是用 $\varepsilon$-约束的混合整数线性规划求解多目标，再以启发式算法作对照。这套流程可直接用于食品、香料、烟草等风味体系的成分—感官建模，也可用于药物辅料配比一类的多目标配方设计。

深度条件生成模型在本文场景下未被采用，原因见第八章：样本量不足会使生成器退化为插值器。若数据规模扩大到每风格数百样本，可把收缩高斯换成条件变分自编码器，优化部分仍保留现有的混合整数线性规划结构。
